% !TEX root = ../../summer_2026_introduction.tex
\subsection{\textbf{注意点と共通点}}
まずは、分析にあたっての注意点を述べる。2021年度までは学校裁量問題制度が導入されていたため、
問題構成や配点が最近の傾向とは異なることに注意してほしい。

次に、英語、数学共に最近の問題構成、配点及び出題分野は \textbf{安定している}。

英語に至っては、配点まで全く同じである。したがって、今年度の入試も同じだと仮定して良い。
仮に、傾向等が変わったとしても、皆平等に混乱するため、結局は実力が試されることに違いはない。

\subsection{\textbf{英語の分析}}
英語の入試について、分析していく。適宜、別冊付録の分析結果を参照してほしい。

問題構成は大問1から4までの4つの大問構成である。

大問1はリスニング。配点は35点。

大問2は語句・文法問題。配点は16点。

大問3はさらにAからCまでの3つの小問に分かれており、それぞれに長文読解問題がある。それぞれの小問の配点は10点、10点、17点と
やや3Cの配点が大きい。

大問4は英作文問題。配点は12点。
\vspace{1em}

分析していて感じたことは、リスニングの配点が大きいことはもちろんのこと、
\textbf{空欄補充及び英作文の問題数が大きと感じた。}英語といえば、「長文読解」というイメージだが、
これらの問題数が多いことは、英語の基本を問う問題が多いことを示している。英語が苦手な生徒と得意な生徒の差がつきにくいのではないかと思う。

一方、長文読解は文量はそこまで多くないにしろ、得手不得手がはっきりと出る問題であるため、差がつきやすい。

\newpage

\subsection{\textbf{数学の分析}}
最後に、数学の入試について、分析していく。適宜、別冊付録の分析結果を参照してほしい。

問題構成は大問1から5までの5つの大問構成である。

大問1は小問集合。配点は35点。

大問2はいろいろ。資料の整理、データの比較の出題が目立つ。配点は15点。

大問3は関数。平面図形、三角形・四角形の複合問題が多い印象。配点は16点。

大問4は幾何。配点は16点。

大問5は幾何。証明問題が必ず出題される。配点は16点。
\vspace{0.5em}

ここで、次のランキングを見てほしい。
\vspace{0.5em}

    \begin{table}[htbp]
        \centering
        \caption{北海道数学：総合ランキング \\ 分野統合版}
        \begin{tabular}{c l r}
        \hline
        \textbf{順位} & \textbf{分野} & \textbf{スコア} \\
        \hline
        1 & 基本図形系（平面図形＋三角形・四角形＋平行と合同） & 50 \\
        2 & 計算系（正負の数＋文字式＋式の計算）& 32 \\
        3 & 関数系（関数 + 1次関数 + 2乗に比例する関数 & 25 \\
        4 & 統計系（資料の整理＋データの比較）& 21 \\
        5 & 平方根 & 15 \\
        6 & 確率 & 12 \\
        7 & 空間図形 & 9 \\
        7 & 相似 & 9 \\
        9 & 連立方程式 & 8 \\
        10 & 二次方程式 & 7 \\
        11 & 円 & 6 \\
        12 & 三平方の定理 & 4 \\
        & その他 & 2 \\
        \hline
        \end{tabular}
    \end{table}


分析していて感じたことは、\textbf{基本図形系の配点が高い}ことである。しかし、このランキングはあくまでも
\textbf{出題頻度} に重点を置いたものであり、実際の配点までは考慮出来ていない。

数学は、他の教科と比べて得点が取りにくい教科である。得意な人間でも8割取れたら大したもんである。
したがって、他と差をつけるために難しい問題を挑戦するよりも、\textbf{基礎的な問題を確実に解く方が極めて重要である。}
\newpage
